\documentclass[11pt]{article}
\input{quasse/shared/preamble}
\title{QuaSSE: mathematical conventions and derivations}
\begin{document}
\maketitle

This selective reference records mathematical reasoning that can be consulted and extended during
development. It derives the trait generator and the root/origin conditioning normalizations; it is not a
complete derivation of the QuaSSE likelihood. The \href{guide.pdf}{guide} gives an introductory account of the model.

\section{Conventions and conditional quantities}\label{sec:conventions}
The latent trait is $x\in\mathbb{R}$. Forward elapsed time is $s$; age before the present is $t$.
Birth and death rates are $\lambda(x)$ and $\mu(x)$. Fossil sampling is non-removing with scalar
rate $\psi\geq0$, and present sampling has scalar probability $\rho\in[0,1]$. Sampling events
and trait observations are distinct: $g(y\mid x)$ denotes the observation density.

Conditional on a single lineage at age $t$ with trait $x$, define $E(x,t)$ as the probability of no
sampled descendants between that time and the present, including no fossil samples. Consequently
$0\leq E\leq1$ and $E(x,0)=1-\rho$. Define $D(x,t)$, for a specified observed descendant subtree,
as its likelihood contribution conditional on that lineage and trait, before any whole-tree root
conditioning. This contribution integrates over unobserved events and trait histories. It is
nonnegative but need not be at most one. Node ordering and whole-tree conditioning conventions
must additionally be specified when deriving a full likelihood; they are not fixed by this local
definition alone.

These meanings and the living-tip boundary $D(x,0)=\rho g(y\mid x)$ follow the conventions described
in the \href{../../examples/QuaSSE_fossils_fixed_tree.md}{package fossil guide}. They concern the
unscaled mathematical quantities: numerical normalization of $D$ does not change its definition,
and $E$ is a probability rather than a freely rescalable likelihood.

\section{Brownian trait generator}\label{sec:generator}
Assume constant drift $v\in\mathbb{R}$ and constant variance rate $q>0$, on the unbounded trait line.
The forward process, conditional on $X_0=x$, is
\begin{equation}\label{eq:trait-process}
 X_s=x+vs+\sqrt{q}\,W_s,
\end{equation}
where $W_s$ is standard Brownian motion. Drift has units trait/time and $q$ has units
trait$^2$/time. In the package, the input named \texttt{diffusion} is $q$.

For a smooth function $f$, define the backward transition operator
$P_s f(x)=\expect[f(X_s)\mid X_0=x]$. For the following elementary derivation, it suffices to take
$f$ three times continuously differentiable with bounded derivatives. For a small interval
$\delta>0$, write $Z=v\delta+\sqrt{q\delta}\,\xi$ with $\xi$ standard normal. Taylor expansion gives
\begin{align*}
 \expect[f(x+Z)]&=f(x)+f'(x)\expect[Z]
                   +\tfrac12 f''(x)\expect[Z^2]+O(\expect[|Z|^3]),\\
 \expect[Z]&=v\delta,\qquad
 \expect[Z^2]=q\delta+v^2\delta^2,\qquad
 \expect[|Z|^3]=O(\delta^{3/2}).
\end{align*}
Subtracting $f(x)$, dividing by $\delta$, and taking the limit proves
\begin{equation}\label{eq:generator}
 \traitop f:=\lim_{\delta\downarrow0}\frac{P_\delta f-f}{\delta}
       =v f_x+\frac q2 f_{xx}.
\end{equation}
This is an infinitesimal identity under the stated assumptions, not a numerical error bound.

\section{Backward likelihood transport and forward densities}
For trait-only propagation of a descendant contribution $f$, let $u(x,t)=P_t f(x)$. Extending the
ancestral-to-descendant interval increases $t$; the Markov property and
Equation~\eqref{eq:generator} give
\begin{equation}\label{eq:backward-transport}
 u_t=\traitop u=v u_x+\frac q2 u_{xx},\qquad u(x,0)=f(x).
\end{equation}
For an already positive-aged descendant, $t$ in this trait-only formula may be replaced by the
elapsed branch duration. The time-homogeneous operator is unchanged. Full likelihood propagation
also includes diversification and sampling terms, which are outside this derivation.

By contrast, the forward density $p(x,s)$ of the evolving trait obeys the adjoint equation
\begin{equation}\label{eq:forward-density}
 p_s=-v p_x+\frac q2 p_{xx}.
\end{equation}
Indeed, integration by parts in $\int p\traitop f\,dx$, with vanishing boundary terms, moves the
first derivative with a minus sign and the second derivative with a plus sign. The different drift
signs reflect different quantities, not conflicting conventions for Brownian motion.

As a check, take $f(x)$ to be a Gaussian profile with center $m$ and variance $a>0$. Then $P_t f$
has center $m-vt$ and variance $a+qt$, with the same integrated weight. A forward Gaussian density
instead moves its center to $m+vs$. The former convention is used in the
\href{../../validation/log_pde/README.md}{log-space transport experiments}.

The unbounded-domain calculation does not justify any particular finite-grid boundary rule.
Spatial discretization, time integration, and boundary truncation require their own analysis.
The original QuaSSE model is described by FitzJohn~\cite{fitzjohn2010}; the derivation here states
its constant-coefficient trait convention explicitly rather than attributing fossil extensions to
that paper.

\section{Sampling success and starting conditions}\label{sec:conditioning}
The quantities $E$ and $D$ retain their meanings in Section~\ref{sec:conventions}. Let $Q(x,t)$
be the probability of no sampled \emph{living} descendants from a lineage at age $t$ and trait $x$.
For non-removing fossil sampling, the relevant backward equations and boundaries are
\begin{align}
 E_t &= \traitop E+\mu-(\lambda+\mu+\psi)E+\lambda E^2,
       & E(x,0)&=1-\rho,\label{eq:E}\\
 Q_t &= \traitop Q+\mu-(\lambda+\mu)Q+\lambda Q^2,
       & Q(x,0)&=1-\rho.\label{eq:Q}
\end{align}
The trait operator is the same $\traitop$ as before. To see why fossil sampling drops out of the
second equation, consider a short forward interval. A fossil observation leaves the lineage alive
and does not alter the event of having no sampled living descendants. Its contribution
$\psi\,\delta Q$ cancels the $-\psi\,\delta Q$ term from the probability of no event. Death
contributes $\mu\,\delta$, and birth contributes $\lambda\,\delta Q^2$ because daughter futures are
conditionally independent. Taking the infinitesimal limit gives Equation~\eqref{eq:Q}.
Thus the existing $E$ solver can evaluate $Q$ with sampling rate zero, without changing the
sampling rate in the data likelihood. The latter still propagates
\begin{equation}
 D_t=\traitop D+[2\lambda E-(\lambda+\mu+\psi)]D.
\end{equation}
The sample-success choices follow the biological distinctions in the BEAST sampled-ancestors
inputs~\cite{saInputs}; SSE retains its historical root/any-sample default.

\subsection{A single-lineage origin}
Let $t_o$ be the origin age, at least the sampled root age $t_r$. Prune the observed tree and
propagate its root partial to $t_o$ to obtain $D_o(x)$. There is no birth or observation event at
the origin. For a normalized, data-independent starting-trait density $\pi_o$, write
\begin{equation}\label{eq:origin-numerator}
 N_o=\int \pi_o(x)D_o(x)\,dx.
\end{equation}
If $S$ is the required sampling-success event and the observed data satisfy $S$, conditional
probability gives $p(\text{data}\mid S)=p(\text{data})/P(S)$. Marginalizing over the trait yields
\begin{equation}\label{eq:origin-condition}
 L_o=\frac{N_o}{Z_o},\qquad
 Z_o=\begin{cases}
 \int\pi_o(x)[1-E(x,t_o)]\,dx & \text{any sampled descendant},\\
 \int\pi_o(x)[1-Q(x,t_o)]\,dx & \text{sampled living descendant},\\
 1 & \text{no sample-success condition}.
 \end{cases}
\end{equation}
An average of pointwise ratios using the original $\pi_o$ would condition differently. Indeed,
after conditioning, the trait density is proportional to $\pi_o(x)P(S\mid x)$; using this density
to average $D_o(x)/P(S\mid x)$ recovers Equation~\eqref{eq:origin-condition}.

Observed events retain their existing contributions: $\rho g$ at a living tip, $\psi gE$ at a
terminal fossil, $\psi gD_{\rm continuing}$ at a sampled ancestor, and
$\lambda D_{\rm left}D_{\rm right}$ at an ordinary bifurcation. Stem fossils are therefore allowed,
including a sampled-ancestor root. These factors also specify which event densities are being
conditioned on; a stem does not add a hidden root event.

\subsection{Starting at a bifurcation}
In SSE's existing convention, the root partial $D_r$ already includes the factor $\lambda$ from
merging the daughters. Retain the weighting density $\pi$ used with this partial. For a normalized,
data-independent $\pi$ and $0<C=\int\pi(x)\lambda(x)\,dx<\infty$, the effective starting-trait
density at the bifurcation is $\pi(x)\lambda(x)/C$. The unconditioned descendant likelihood at
that starting split is consequently $N_r/C$, where $N_r=\int\pi D_r\,dx$.

Both daughters must satisfy the selected success condition. Conditional on their common starting
trait, their futures are independent, so their success probabilities multiply. The probability of
success under the effective trait density is $Z_r/C$, giving
\begin{equation}\label{eq:root-condition}
 L_r=\frac{N_r}{Z_r},\qquad
 Z_r=\begin{cases}
 \int\pi(x)\lambda(x)[1-E(x,t_r)]^2\,dx & \text{any sample on each side},\\
 \int\pi(x)\lambda(x)[1-Q(x,t_r)]^2\,dx & \text{living sample on each side},\\
 \int\pi(x)\lambda(x)\,dx & \text{no sample-success condition}.
 \end{cases}
\end{equation}
The common normalization of $\pi$ cancels in these ratios. This explains why retaining the legacy
root quadrature convention does not require renormalizing its Flat weights. Root/living conditioning
assigns zero support to a tree with an all-fossil daughter, even when the denominator is positive.
A sampled-ancestor root is not a starting bifurcation and is excluded in root mode.

For constant $\lambda$, its factor cancels between numerator and denominator, yielding the usual
constant-rate root expression. With trait dependence, this convention does not identify $\pi$
itself as the trait density conditional on a starting split. Removing sample-success conditioning
still conditions on that split. In particular, this option differs from the raw birth-density
calculation obtained by disabling survival conditioning in diversitree.

When $t_o=t_r$, origin propagation is the identity and preserves all observed root factors. It
does not change the single-lineage starting model into a bifurcation-conditioned model; their
normalizations and support requirements still differ.

\subsection{Trait weights and numerical scales}
At the origin, useful-bin midpoint quadrature with spacing $\Delta x$ normalizes the Flat density
as $\pi_i=1/(n\Delta x)$. The compatibility option Observed instead uses
\begin{equation}\label{eq:observed-weights}
 \pi_i=\frac{D_i}{\sum_j D_j\Delta x}.
\end{equation}
Weights are formed after stem propagation at the endpoint resolution and used unchanged for both
integrals. A common positive scale in $D$ cancels in Equation~\eqref{eq:observed-weights}. However,
because these weights depend on the data, their substitution defines a comparison score, not the
conditional likelihood under an independently specified generative prior.

Flat refers to the finite useful trait interval, not a proper uniform prior on the real line.
Automatic grid bounds depend on observations, and padding can change useful-bin boundaries.
Thus Flat on an automatically selected interval should not be described as a wholly
data-independent scientific prior. Fixed-support comparisons and numerical domain checks are
needed to separate weighting conventions from discretization effects.

If the final numerical partial is $\widetilde D$ and accumulated normalization is $e^c$, the
reported log likelihood is $c+\log\int\pi\widetilde D\,dx-\log Z$. Neither $E$ nor $Q$ is
rescaled as a likelihood. Their numerical calculation, like stem transport, uses the existing
spatial and temporal approximations; the conditional-probability identities above do not supply
a global numerical error bound.

\bibliographystyle{plain}
\bibliography{references}
\end{document}
